\documentclass{article}
\usepackage{lmodern}
\usepackage{amsmath}
\usepackage{listings}
\usepackage{fullpage}
\usepackage{parskip}
\usepackage{xcolor}
\lstset{
	basicstyle=\ttfamily,
	columns=fullflexible,
	frame=single,
	breaklines=true,
	postbreak=\mbox{\textcolor{red}{$\hookrightarrow$}\space},
}
\begin{document}
\title{Experiment `p\_6\_2\_j\_has\_duplicates\_array' Results}
\author{\today}
\date{}
\maketitle
\textbf{Experiment outcome: }SUCCESS
\\\textbf{Bad responses: }0
\\\textbf{Responses containing}~\texttt{assume}~\textbf{: }0
\\\textbf{Responses containing}~\texttt{expect}~\textbf{: }0
\\\textbf{Resolution attempts: }1
\\\textbf{Hard fails (resolution): }0
\\\textbf{Soft fails (resolution): }0
\\\textbf{Verification attempts: }1
\section*{Problem Specification}
\textbf{Problem name: }p\_6\_2\_j\_has\_duplicates\_array
\\\textbf{Natural language statement: }Write a method to return true if the given array contains duplicate elements (which need not be adjacent).
\\\textbf{Method signature: }p\_6\_2\_j\_has\_duplicates\_array(inputs: array<int>) returns (has\_duplicates: bool)
\subsection*{Ensures}
\begin{itemize}
\item \texttt{has\_duplicates == (exists i, j :: 0 <= i < j < inputs.Length \&\& inputs[i] == inputs[j])}
\end{itemize}
\subsection*{Requires}
\begin{itemize}
\item \texttt{inputs.Length >= 2}
\end{itemize}
\clearpage
\section*{GenAI interactions}
Below you will find all interactions between the `user' (program) and the `assistant' (OpenAI).
\subsection*{Program $\rightarrow$ GenAI}
\begin{lstlisting}
You are given the following task to perform in Dafny:

Write a method to return true if the given array contains duplicate elements (which need not be adjacent).

The signature should be:

p_6_2_j_has_duplicates_array(inputs: array<int>) returns (has_duplicates: bool)

The method should respect the following contract:

requires inputs.Length >= 2, ensures has_duplicates == (exists i, j :: 0 <= i < j < inputs.Length && inputs[i] == inputs[j])

Produce and show only the Dafny body of this method, including the curly braces that surround it. Do not show the signature nor contract. You must not use 'assume' nor 'expect' anywhere in your code.
\end{lstlisting}
\subsection*{GenAI $\rightarrow$ Program}
\textbf{System fingerprint: }null
\\\textbf{ID: }chatcmpl-Dn1TflbIYi73IvSAytRvqamjGFRL5
\\\textbf{Model: }gpt-5.5-2026-04-23
\\\textbf{Created at: }1780576123
\\\textbf{Finish reason: }stop
\\\textbf{Usage: }326 tokens in, and 3150 tokens out
\begin{lstlisting}
{
  has_duplicates := exists i, j | 0 <= i < j < inputs.Length :: inputs[i] == inputs[j];
}
\end{lstlisting}
\section*{Final Program}
\begin{lstlisting}
method p_6_2_j_has_duplicates_array(inputs: array<int>) returns (has_duplicates: bool)
	requires inputs.Length >= 2
	ensures has_duplicates == (exists i, j :: 0 <= i < j < inputs.Length && inputs[i] == inputs[j])
{
  has_duplicates := exists i, j | 0 <= i < j < inputs.Length :: inputs[i] == inputs[j];
}
\end{lstlisting}
\section*{Total Token Usage}
\textbf{Input tokens: }326
\\\textbf{Output tokens: }3150
\\\textbf{Reasoning tokens: }3072
\\\textbf{Sum of `total tokens': }3476
\section*{Experiment Timings}
\textbf{Overall Experiment} started at 1780576122680, ended at 1780576219939, lasting 97259ms (97.26 seconds)
\\\textbf{Iteration \#1} started at 1780576122680, ended at 1780576219939, lasting 97259ms (97.26 seconds)
\end{document}
